\section{Background: ABF for a reaction-coordinate free energy}
\label{sec:abf}

\subsection{Free energy and the unbiased marginal}
With reaction coordinate $\xib(x,y)=x$, the unbiased reaction-coordinate
marginal of the canonical density~\eqref{eq:canonical} is
\begin{equation}
  \pref(x) \;=\;
  \frac{\displaystyle\int e^{-\beta V(x,y)}\dd y}
       {\displaystyle\iint e^{-\beta V(x,y)}\dd x\dd y},
  \label{eq:pref}
\end{equation}
and the free energy along $\xib$ is, up to an additive constant,
\begin{equation}
  \Fref(x) \;=\; -\frac{1}{\beta}\,\log \pref(x) + C.
  \label{eq:Fref}
\end{equation}
Because $F$ is defined only up to the constant $C$, every free-energy profile in
this report --- both estimates and the reference --- is \emph{centred} (its mean
over the interior evaluation window is subtracted) before any $\Ltwo$ error is
formed; otherwise a meaningless constant offset would dominate the comparison.
The reference quantities are computed by a one-dimensional $y$-quadrature on a
fine grid with a log-sum-exp shift for numerical stability, which is exact up to
discretisation and provides the ground truth $\Fref,\Fpref,\pref$.

\subsection{The mean force is the object ABF estimates}
Differentiating~\eqref{eq:Fref} and using $\partial_x \pref$ gives the
\emph{mean force}
\begin{equation}
  \Fpref(x) \;=\;
  \frac{\displaystyle\int \partial_x V(x,y)\,e^{-\beta V(x,y)}\dd y}
       {\displaystyle\int e^{-\beta V(x,y)}\dd y}
  \;=\; \E_{\pi}\!\left[\,\partial_x V(X,Y)\,\middle|\, X=x\,\right].
  \label{eq:meanforce}
\end{equation}
Equation~\eqref{eq:meanforce} is the cornerstone of ABF: the gradient of the
free energy equals the \emph{conditional expectation of the instantaneous force}
$\partial_x V$ given the reaction coordinate. It is a local, sample-based
quantity --- no integral of $V$ is required --- so it can be estimated on the
fly from the particles currently at $X\approx x$, and the free energy is then
recovered by quadrature, $\Fhat(x)=\int^x \Fphat$.

\subsection{Biased dynamics}
ABF simulates $N$ overdamped Langevin replicas that share one running mean-force
estimate $\Fphat_t(x)\approx\Fpref(x)$ with bias potential
$\Bt(x)=\int^x\Fphat_t$. Each replica $(X_t,Y_t)$ evolves as
\begin{equation}
  \begin{aligned}
    \dd X_t &= \bigl[-\partial_x V(X_t,Y_t) + \Fphat_t(X_t)\bigr]\dd t
              + \sqrt{2/\beta}\,\dd W_t^x, \\
    \dd Y_t &= -\,\partial_y V(X_t,Y_t)\,\dd t
              + \sqrt{2/\beta}\,\dd W_t^y,
  \end{aligned}
  \label{eq:abf_sde}
\end{equation}
with independent Brownian motions $W^x,W^y$. The added force $+\Fphat_t(X_t)$
acts \emph{only} in the $x$-direction; the $y$-dynamics is unbiased. When
$\Fphat_t\to\Fpref$ the net drift along $x$ becomes $-\partial_x(V-F)$, whose
$x$-marginal equilibrium is flat, so the reaction coordinate diffuses freely
across the barrier.

\subsection{Why the bias does not spoil the mean-force estimate}
A bias that depends only on $x$ leaves the conditional law of $Y$ given $X=x$
invariant. This is what makes ABF self-consistent.

\begin{proposition}[Conditional invariance under an $x$-only bias]
\label{prop:cond_invariance}
Let $\nu$ be the canonical measure $\nu\propto e^{-\beta V(x,y)}$ and let
$\nu_B\propto e^{-\beta\,[V(x,y)-B(x)]}$ be the measure biased by any function
$B$ of $x$ alone. Then for every $x$ the conditionals coincide,
$\nu_B(y\mid X=x)=\nu(y\mid X=x)$, and consequently
\begin{equation}
  \E_{\nu_B}\!\left[\partial_x V(X,Y)\mid X=x\right]
  \;=\;
  \E_{\nu}\!\left[\partial_x V(X,Y)\mid X=x\right]
  \;=\; \Fpref(x).
  \label{eq:cond_equal}
\end{equation}
\end{proposition}

\begin{proof}
At fixed $X=x$ the bias factor $e^{\beta B(x)}$ is a constant in $y$ and cancels
between numerator and normaliser of the conditional density:
\[
  \nu_B(y\mid X=x)
  = \frac{e^{-\beta[V(x,y)-B(x)]}}{\int e^{-\beta[V(x,y')-B(x)]}\dd y'}
  = \frac{e^{-\beta V(x,y)}}{\int e^{-\beta V(x,y')}\dd y'}
  = \nu(y\mid X=x).
\]
Taking the conditional expectation of $\partial_x V$ under this common
conditional gives~\eqref{eq:cond_equal}, the last equality
being~\eqref{eq:meanforce}.
\end{proof}

\Cref{prop:cond_invariance} is the reason ABF works: even though the biased
ensemble~\eqref{eq:abf_sde} has a different $x$-marginal than the canonical
one, the conditional samples $Y\mid X=x$ remain canonical, so the empirical
conditional average of $\partial_x V$ is an unbiased target for
$\Fpref(x)$. Any modification of ABF must respect this property; in particular,
a birth--death correction that resamples whole configurations $(x,y)$ is
only admissible if it too leaves $Y\mid X=x$ (approximately) intact ---
a hypothesis we test directly in each of the three case studies
(\cref{sec:case_meta,sec:case_eb,sec:case_wca}).

\begin{remark}[Generalisation to many-body systems]
\label{rem:manybody}
Two of our three systems are two-dimensional with $\xib=x$, but the WCA dimer of
\cref{sec:case_wca} is a many-body system: a configuration is the full vector
$q\in\R^{2n}$ of all particle coordinates, and the reaction coordinate is a
scalar \emph{collective variable} $z(q)$ (a scaled bond length) rather than a
single Cartesian component. \Cref{prop:cond_invariance} holds verbatim under the
substitution $x\mapsto z(q)$ and $y\mapsto$ the orthogonal complement of $z$ on
the constant-$z$ manifold (here the solvent coordinates and the transverse dimer
modes): a bias $B(z)$ that depends on the collective variable alone is constant
on each level set of $z$, so it cancels in the conditional density and leaves the
law of all orthogonal degrees of freedom invariant. Consequently the mean force
along $z$ --- which for a nonlinear $z(q)$ also carries a geometric term
(\cref{sec:case_wca}) --- is still estimated without bias. The only requirement
the FR step must meet is the same one as in two dimensions: a clone must copy the
\emph{entire} configuration $q\in\R^{2n}$, so that the orthogonal coordinates
travel with $z$ and the conditional law is preserved.
\end{remark}

\subsection{The ABF mean-force estimator}
\label{sec:abf_estimator}
In the kernel form used by the reference implementation, $\Fphat$ is a
Nadaraya--Watson running average of the instantaneous force,
\begin{equation}
  \Fphat_n(x) \;=\;
  \frac{\sum_{k\le n}\sum_i K_h(x-X_k^i)\,\partial_x V(X_k^i,Y_k^i)}
       {\sum_{k\le n}\sum_i K_h(x-X_k^i)},
  \label{eq:nw}
\end{equation}
with a Gaussian kernel $K_h$ of bandwidth $h$, accumulated over all replicas
$i$ and past steps $k$; the free-energy estimate is its pinned trapezoidal
antiderivative, $\Fhat_n(x)=\int_0^x\Fphat_n$, and we use $\Bt\equiv\Fhat$ as
the bias. The GPU backend used for the production study replaces the
$O(n_{\mathrm{grid}}\,N)$ kernel sums by an $O(N+n_{\mathrm{grid}})$
\emph{binned-smoothed} estimator: particles are scattered onto the $x$-grid and
the count and force histograms are Gaussian-smoothed, so that
$\Fphat(x)=\mathrm{smooth(force)}/\mathrm{smooth(count)}$ is a Riemann-sum
approximation of~\eqref{eq:nw} that converges to it as the grid spacing
vanishes. The two estimators are validated to agree to well under $1\%$ on
static tests (\cref{sec:design}).
